\documentclass[10pt,a4paper]{amsart}
\usepackage{amsmath,amssymb,amsthm,graphicx}
\usepackage{verbatim}
\newenvironment{reference}{\comment}{\endcomment}
%\newenvironment{reference}{}{}
\newenvironment{slogan}{\comment}{\endcomment}
\newenvironment{history}{\comment}{\endcomment}

% For commutative diagrams we use Xy-pic
\usepackage[all]{xy}

% We use 2cell for 2-commutative diagrams.
\xyoption{2cell}
% Theorem environments.
%
\theoremstyle{plain}
\newtheorem{theorem}[subsection]{Theorem}
\newtheorem{proposition}[subsection]{Proposition}
\newtheorem{lemma}[subsection]{Lemma}

\theoremstyle{definition}
\newtheorem{definition}[subsection]{Definition}
\newtheorem{example}[subsection]{Example}
\newtheorem{exercise}[subsection]{Exercise}
\newtheorem{situation}[subsection]{Situation}

\theoremstyle{remark}
\newtheorem{remark}[subsection]{Remark}
\newtheorem{remarks}[subsection]{Remarks}

\numberwithin{equation}{subsection}

% Macros
%
\def\lim{\mathop{\mathrm{lim}}\nolimits}
\def\colim{\mathop{\mathrm{colim}}\nolimits}
\def\Spec{\mathop{\mathrm{Spec}}}
\def\Hom{\mathop{\mathrm{Hom}}\nolimits}
\def\Ext{\mathop{\mathrm{Ext}}\nolimits}
\def\SheafHom{\mathop{\mathcal{H}\!\mathit{om}}\nolimits}
\def\SheafExt{\mathop{\mathcal{E}\!\mathit{xt}}\nolimits}
\def\Sch{\mathit{Sch}}
\def\Mor{\mathop{\mathrm{Mor}}\nolimits}
\def\Ob{\mathop{\mathrm{Ob}}\nolimits}
\def\Sh{\mathop{\mathit{Sh}}\nolimits}
\def\NL{\mathop{N\!L}\nolimits}
\def\CH{\mathop{\mathrm{CH}}\nolimits}
\def\proetale{{pro\text{-}\acute{e}tale}}
\def\etale{{\acute{e}tale}}
\def\QCoh{\mathit{QCoh}}
\def\Ker{\mathop{\mathrm{Ker}}}
\def\Im{\mathop{\mathrm{Im}}}
\def\Coker{\mathop{\mathrm{Coker}}}
\def\Coim{\mathop{\mathrm{Coim}}}

% Boxtimes
%
\DeclareMathSymbol{\boxtimes}{\mathbin}{AMSa}{"02}

%
% Macros for moduli stacks/spaces
%
\def\QCohstack{\mathcal{QC}\!\mathit{oh}}
\def\Cohstack{\mathcal{C}\!\mathit{oh}}
\def\Spacesstack{\mathcal{S}\!\mathit{paces}}
\def\Quotfunctor{\mathrm{Quot}}
\def\Hilbfunctor{\mathrm{Hilb}}
\def\Curvesstack{\mathcal{C}\!\mathit{urves}}
\def\Polarizedstack{\mathcal{P}\!\mathit{olarized}}
\def\Complexesstack{\mathcal{C}\!\mathit{omplexes}}
% \Pic is the operator that assigns to X its picard group, usage \Pic(X)
% \Picardstack_{X/B} denotes the Picard stack of X over B
% \Picardfunctor_{X/B} denotes the Picard functor of X over B
\def\Pic{\mathop{\mathrm{Pic}}\nolimits}
\def\Picardstack{\mathcal{P}\!\mathit{ic}}
\def\Picardfunctor{\mathrm{Pic}}
\def\Deformationcategory{\mathcal{D}\!\mathit{ef}}
\usepackage{cleveref}
\crefname{originallemma}{original lemma}{original lemmas}
\Crefname{originallemma}{Original lemma}{Original lemmas}
\makeatletter
\newlabel{descent-lemma-descending-property-flat}{{02L2}{source}{}{sourcepointer-1}{}}
\newlabel{descent-lemma-descending-property-flat@cref}{{[originallemma][0][]02L2}{[1][0][]source}{}{}{}}
\makeatother
\crefname{originallemma}{original lemma}{original lemmas}
\Crefname{originallemma}{Original lemma}{Original lemmas}
\makeatletter
\newlabel{descent-lemma-descending-property-locally-finite-presentation}{{02KY}{source}{}{sourcepointer-2}{}}
\newlabel{descent-lemma-descending-property-locally-finite-presentation@cref}{{[originallemma][0][]02KY}{[1][0][]source}{}{}{}}
\makeatother
\crefname{originallemma}{original lemma}{original lemmas}
\Crefname{originallemma}{Original lemma}{Original lemmas}
\makeatletter
\newlabel{descent-lemma-descending-property-smooth}{{02VL}{source}{}{sourcepointer-3}{}}
\newlabel{descent-lemma-descending-property-smooth@cref}{{[originallemma][0][]02VL}{[1][0][]source}{}{}{}}
\makeatother
\crefname{originalreference}{original reference}{original references}
\Crefname{originalreference}{Original reference}{Original references}
\makeatletter
\newlabel{lemma-base-change-nodal-family}{{0C5B}{source}{}{sourcepointer-4}{}}
\newlabel{lemma-base-change-nodal-family@cref}{{[originalreference][0][]0C5B}{[1][0][]source}{}{}{}}
\makeatother
\crefname{originalreference}{original reference}{original references}
\Crefname{originalreference}{Original reference}{Original references}
\makeatletter
\newlabel{lemma-node-field-extension}{{0C56}{source}{}{sourcepointer-5}{}}
\newlabel{lemma-node-field-extension@cref}{{[originalreference][0][]0C56}{[1][0][]source}{}{}{}}
\makeatother
\crefname{originallemma}{original lemma}{original lemmas}
\Crefname{originallemma}{Original lemma}{Original lemmas}
\makeatletter
\newlabel{morphisms-lemma-dimension-fibre-after-base-change}{{02FY}{source}{}{sourcepointer-6}{}}
\newlabel{morphisms-lemma-dimension-fibre-after-base-change@cref}{{[originallemma][0][]02FY}{[1][0][]source}{}{}{}}
\makeatother
\title{Being at-worst-nodal of relative dimension 1 is fpqc local on the target}
\author{The Stacks Project Authors}
\date{stacks-project:curves.tex:local-snapshot-20261012; GNU FDL 1.2 or later}
\begin{document}
\maketitle
\noindent\textit{Editorial note.} All delivered mathematical text is the author's own, copied from the pinned Stacks Project chapter file sources/stacks-curves/pinned-original-source.tex; no statement or proof step was rewritten, reordered or removed.\par\medskip
\noindent\textit{Editorial note.} The counted proof is the author's complete argument: the field case of the proof invokes \textbackslash{}label\{lemma-node-field-extension\} (source line 5090), a published lemma of the same chapter whose own proof delegates its central ring-isomorphism step to the reader (source lines 4646-4652: 'We encourage the reader to prove this for themselves ... We omit the details.'); that lemma was therefore not carded separately, but it is used here as published input per the citation policy; the remaining steps of the present proof are complete.\par\medskip
\noindent\textit{Editorial note.} The label lemma-nodal-family-fpqc-local-target names the counted statement, and the proof environment immediately following it in the source is delivered verbatim as its proof.\par\medskip
\noindent\textit{Editorial note.} The counted statement is carried by the author's own lemma environment in the chapter source.\par\medskip
\noindent\textit{Allowed standard prerequisites (editorial).}\begin{itemize}
\item Base change of nodal families. Same chapter, Lemma \textbackslash{}label\{lemma-base-change-nodal-family\}. Proof lines 5078-5079: the direction (1) implies (2) follows from the base change lemma for nodal families.
\item Descent of the structural properties along fpqc coverings. Descent, Lemmas \textbackslash{}label\{descent-lemma-descending-property-locally-finite-presentation\} and \textbackslash{}label\{descent-lemma-descending-property-flat\}; Morphisms, Lemma \textbackslash{}label\{morphisms-lemma-dimension-fibre-after-base-change\}. Proof lines 5081-5087: being locally of finite presentation, flat, or of relative dimension 1 is fpqc local on the target, which reduces the remaining direction to the case of a field.
\item Nodes under field extension. Same chapter, Lemma \textbackslash{}label\{lemma-node-field-extension\}. Proof lines 5090-5093: over a field the statement follows from this lemma, together with the fpqc-locality of smoothness (Descent, Lemma \textbackslash{}label\{descent-lemma-descending-property-smooth\}).
\end{itemize}\medskip
\ifdefined\hypertarget\hypertarget{sourcepointer-1}{}\fi
\noindent\textit{Source reference.} Original lemma 02L2: sources/stacks-chow/tags.txt, source lines 3393--3393. Prerequisite pointer invoked at proof line(s) 5086; it names a standard fact of a sibling chapter that the counted argument uses, and every step of the counted proof built on it is retained verbatim in the delivered body.\par\medskip
\ifdefined\hypertarget\hypertarget{sourcepointer-2}{}\fi
\noindent\textit{Source reference.} Original lemma 02KY: sources/stacks-chow/tags.txt, source lines 3389--3389. Prerequisite pointer invoked at proof line(s) 5085; it names a standard fact of a sibling chapter that the counted argument uses, and every step of the counted proof built on it is retained verbatim in the delivered body.\par\medskip
\ifdefined\hypertarget\hypertarget{sourcepointer-3}{}\fi
\noindent\textit{Source reference.} Original lemma 02VL: sources/stacks-chow/tags.txt, source lines 3772--3772. Prerequisite pointer invoked at proof line(s) 5092; it names a standard fact of a sibling chapter that the counted argument uses, and every step of the counted proof built on it is retained verbatim in the delivered body.\par\medskip
\ifdefined\hypertarget\hypertarget{sourcepointer-4}{}\fi
\noindent\textit{Source reference.} Original reference 0C5B: sources/stacks-curves/pinned-original-source.tex, source lines 4932--4932. Prerequisite pointer invoked at proof line(s) 5080; it names a standard fact of the same chapter that the counted argument uses, and every step of the counted proof built on it is retained verbatim in the delivered body.\par\medskip
\ifdefined\hypertarget\hypertarget{sourcepointer-5}{}\fi
\noindent\textit{Source reference.} Original reference 0C56: sources/stacks-curves/pinned-original-source.tex, source lines 4605--4605. Prerequisite pointer invoked at proof line(s) 5090; it names a standard fact of the same chapter that the counted argument uses, and every step of the counted proof built on it is retained verbatim in the delivered body.\par\medskip
\ifdefined\hypertarget\hypertarget{sourcepointer-6}{}\fi
\noindent\textit{Source reference.} Original lemma 02FY: sources/stacks-chow/tags.txt, source lines 3209--3209. Prerequisite pointer invoked at proof line(s) 5087; it names a standard fact of a sibling chapter that the counted argument uses, and every step of the counted proof built on it is retained verbatim in the delivered body.\par\medskip
\begin{lemma}
\label{lemma-nodal-family-fpqc-local-target}
Let $f : X \to S$ be a morphism of schemes. Let $\{U_i \to S\}$
be an fpqc covering. The following are equivalent
\begin{enumerate}
\item $f$ is at-worst-nodal of relative dimension $1$,
\item each $X \times_S U_i \to U_i$ is at-worst-nodal of relative
dimension $1$.
\end{enumerate}
In other words, being at-worst-nodal of relative dimension $1$
is fpqc local on the target.
\end{lemma}

\begin{proof}
One direction we have seen in Lemma \ref{lemma-base-change-nodal-family}.
For the other direction, observe that being locally of finite
presentation, flat, or to have relative dimension $1$
is fpqc local on the target
(Descent, Lemmas
\ref{descent-lemma-descending-property-locally-finite-presentation},
\ref{descent-lemma-descending-property-flat}, and
Morphisms, Lemma \ref{morphisms-lemma-dimension-fibre-after-base-change}).
Taking fibres we reduce
to the case where $S$ is the spectrum of a field. In this case the
result follows from Lemma \ref{lemma-node-field-extension}
(and the fact that being smooth is fpqc local on the target by
Descent, Lemma \ref{descent-lemma-descending-property-smooth}).
\end{proof}

\end{document}
